% GENERATED FILE. Edit canonical lessons, then run npm run generate:books.
% canonical-source-hash: fnv1a64:92c2209b90aa9dd5
% canonical-lessons: SA-C444-purnima, SA-C444-amavasya, SA-C444-suklapaksah, SA-C444-krsnapaksah, SA-C444-rajate

\chapter{The Moon's Fortnights}
\label{ch:sa-the-moon-s-fortnights}
\hlchaptermodality{444}

\begin{chapteropening}
\textbf{By the end of this chapter:} \emph{I can say what the moon is doing this fortnight.}

Complete the last lesson of chapter 444: I can say what the moon is doing this fortnight.
\end{chapteropening}

\section[\texorpdfstring{\sk{पूर्णिमा} (pūrṇimā)}{पूर्णिमा (pūrṇimā)}]{\sk{पूर्णिमा} (pūrṇimā) — the full-moon day}

\label{lesson:SA-C444-purnima}



\begin{quote}
Before the new one: say the Sanskrit for a dark night, then the Sanskrit for by night.
\end{quote}

\subsection*{You'll want to know: \sk{पूर्णिमा}}
\textbf{\sk{पूर्णिमा}} — \emph{pūrṇimā} — \textquotedblleft{}the full-moon day\textquotedblright{}.

\begin{quote}
\sk{पूर्णिमायां} \sk{चन्द्रः} \sk{पूर्णः} \sk{भवति।}
\end{quote}

\begin{quote}
\emph{pūrṇimāyāṃ candraḥ pūrṇaḥ bhavati.}
\end{quote}

\begin{quote}
On the full-moon day the moon is full.
\end{quote}

A feminine noun.

\begin{cousinweb}
\textbf{\sk{पूर्णिमा}} (\emph{pūrṇimā}) is built on \emph{pūrṇa}, \textquotedblleft{}full\textquotedblright{}; MW gives \emph{pūrṇimā} as the day of full moon (MW).
\end{cousinweb}

\begin{grammarlens}[title={What you've built}]
One more word for getting around.
\end{grammarlens}

\subsection*{Your turn}
\begin{itemize}
\raggedright
  \item \emph{Say it:} \emph{pūrṇimā}
  \item \emph{Say it:} \emph{pūrṇimā}, once more
  \item \emph{Say it:} \emph{pūrṇimāyāṃ candraḥ pūrṇaḥ bhavati.} — On the full-moon day the moon is full.
  \item \emph{Recall:} say the Sanskrit for a dark night, then the Sanskrit for by night, then say \emph{pūrṇimā} again
\end{itemize}

\begin{tcolorbox}[breakable,colback=teal!4,colframe=teal!35!black,title={Before you move on}]
What does \sk{पूर्णिमा} mean? (\textquotedblleft{}The full-moon day\textquotedblright{}.) Say it once more.
\end{tcolorbox}
\section[\texorpdfstring{\sk{अमावास्या} (amāvāsyā)}{अमावास्या (amāvāsyā)}]{\sk{अमावास्या} (amāvāsyā) — the new-moon night}

\label{lesson:SA-C444-amavasya}



\begin{quote}
Before the new one: say the Sanskrit for by night, then the Sanskrit for the full-moon day.
\end{quote}

\subsection*{You'll want to know: \sk{अमावास्या}}
\textbf{\sk{अमावास्या}} — \emph{amāvāsyā} — \textquotedblleft{}the new-moon night\textquotedblright{}.

\begin{quote}
\sk{अमावास्यायाम्} \sk{आकाशे} \sk{चन्द्रः} \sk{न} \sk{अस्ति।}
\end{quote}

\begin{quote}
\emph{amāvāsyāyām ākāśe candraḥ na asti.}
\end{quote}

\begin{quote}
On the new-moon night there is no moon in the sky.
\end{quote}

A feminine noun.

\begin{cousinweb}
\textbf{\sk{अमावास्या}} (\emph{amāvāsyā}) is built from \emph{amā}, \textquotedblleft{}at home, together\textquotedblright{}, and the root √vas, \textquotedblleft{}to dwell\textquotedblright{}: the night when sun and moon dwell together; MW gives \emph{amāvāsyā} as the night of new moon; it lives on as Hindi \emph{amāwas}, \textquotedblleft{}the night of the new moon\textquotedblright{} (MW; CDIAL 565).
\end{cousinweb}

\begin{grammarlens}[title={What you've built}]
One more word for getting around.
\end{grammarlens}

\subsection*{Your turn}
\begin{itemize}
\raggedright
  \item \emph{Say it:} \emph{amāvāsyā}
  \item \emph{Say it:} \emph{amāvāsyā}, once more
  \item \emph{Say it:} \emph{amāvāsyāyām ākāśe candraḥ na asti.} — On the new-moon night there is no moon in the sky.
  \item \emph{Recall:} say the Sanskrit for by night, then the Sanskrit for the full-moon day, then say \emph{amāvāsyā} again
\end{itemize}

\begin{tcolorbox}[breakable,colback=teal!4,colframe=teal!35!black,title={Before you move on}]
What does \sk{अमावास्या} mean? (\textquotedblleft{}The new-moon night\textquotedblright{}.) Say it once more.
\end{tcolorbox}
\section[\texorpdfstring{\sk{शुक्लपक्षः} (śuklapakṣaḥ)}{शुक्लपक्षः (śuklapakṣaḥ)}]{\sk{शुक्लपक्षः} (śuklapakṣaḥ) — the bright fortnight}

\label{lesson:SA-C444-suklapaksah}



\begin{quote}
Before the new one: say the Sanskrit for the full-moon day, then the Sanskrit for the new-moon night.
\end{quote}

\subsection*{You'll want to know: \sk{शुक्लपक्षः}}
\textbf{\sk{शुक्लपक्षः}} — \emph{śuklapakṣaḥ} — \textquotedblleft{}the bright fortnight\textquotedblright{}.

\begin{quote}
\sk{शुक्लपक्षे} \sk{चन्द्रः} \sk{प्रतिदिनं} \sk{वर्धते।}
\end{quote}

\begin{quote}
\emph{śuklapakṣe candraḥ pratidinaṃ vardhate.}
\end{quote}

\begin{quote}
In the bright fortnight the moon grows every day.
\end{quote}

A masculine noun.

\begin{cousinweb}
\textbf{\sk{शुक्लपक्षः}} (\emph{śuklapakṣaḥ}) is built from \emph{śukla}, \textquotedblleft{}bright, white\textquotedblright{}, and \emph{pakṣa}, \textquotedblleft{}a wing, a side\textquotedblright{}, which MW also gives as the half of a lunar month; MW gives the compound as the light half of the month (MW).
\end{cousinweb}

\begin{grammarlens}[title={What you've built}]
One more word for getting around.
\end{grammarlens}

\subsection*{Your turn}
\begin{itemize}
\raggedright
  \item \emph{Say it:} \emph{śuklapakṣaḥ}
  \item \emph{Say it:} \emph{śuklapakṣaḥ}, once more
  \item \emph{Say it:} \emph{śuklapakṣe candraḥ pratidinaṃ vardhate.} — In the bright fortnight the moon grows every day.
  \item \emph{Recall:} say the Sanskrit for the full-moon day, then the Sanskrit for the new-moon night, then say \emph{śuklapakṣaḥ} again
\end{itemize}

\begin{tcolorbox}[breakable,colback=teal!4,colframe=teal!35!black,title={Before you move on}]
What does \sk{शुक्लपक्षः} mean? (\textquotedblleft{}The bright fortnight\textquotedblright{}.) Say it once more.
\end{tcolorbox}
\section[\texorpdfstring{\sk{कृष्णपक्षः} (kṛṣṇapakṣaḥ)}{कृष्णपक्षः (kṛṣṇapakṣaḥ)}]{\sk{कृष्णपक्षः} (kṛṣṇapakṣaḥ) — the dark fortnight}

\label{lesson:SA-C444-krsnapaksah}



\begin{quote}
Before the new one: say the Sanskrit for the new-moon night, then the Sanskrit for the bright fortnight.
\end{quote}

\subsection*{You'll want to know: \sk{कृष्णपक्षः}}
\textbf{\sk{कृष्णपक्षः}} — \emph{kṛṣṇapakṣaḥ} — \textquotedblleft{}the dark fortnight\textquotedblright{}.

\begin{quote}
\sk{कृष्णपक्षे} \sk{रात्रौ} \sk{अन्धकारः} \sk{भवति।}
\end{quote}

\begin{quote}
\emph{kṛṣṇapakṣe rātrau andhakāraḥ bhavati.}
\end{quote}

\begin{quote}
In the dark fortnight there is darkness at night.
\end{quote}

A masculine noun.

\begin{cousinweb}
\textbf{\sk{कृष्णपक्षः}} (\emph{kṛṣṇapakṣaḥ}) is built from \emph{kṛṣṇa}, \textquotedblleft{}dark, black\textquotedblright{}, and \emph{pakṣa}, \textquotedblleft{}a wing, a side, a half-month\textquotedblright{}; MW gives the compound as the dark half of the month, when the moon wanes (MW).
\end{cousinweb}

\begin{grammarlens}[title={What you've built}]
One more word for getting around.
\end{grammarlens}

\subsection*{Your turn}
\begin{itemize}
\raggedright
  \item \emph{Say it:} \emph{kṛṣṇapakṣaḥ}
  \item \emph{Say it:} \emph{kṛṣṇapakṣaḥ}, once more
  \item \emph{Say it:} \emph{kṛṣṇapakṣe rātrau andhakāraḥ bhavati.} — In the dark fortnight there is darkness at night.
  \item \emph{Recall:} say the Sanskrit for the new-moon night, then the Sanskrit for the bright fortnight, then say \emph{kṛṣṇapakṣaḥ} again
\end{itemize}

\begin{tcolorbox}[breakable,colback=teal!4,colframe=teal!35!black,title={Before you move on}]
What does \sk{कृष्णपक्षः} mean? (\textquotedblleft{}The dark fortnight\textquotedblright{}.) Say it once more.
\end{tcolorbox}
\section[\texorpdfstring{\sk{राजते} (rājate)}{राजते (rājate)}]{\sk{राजते} (rājate) — gleams; shines brightly}

\label{lesson:SA-C444-rajate}



\begin{quote}
Before the new one: say the Sanskrit for the bright fortnight, then the Sanskrit for the dark fortnight.
\end{quote}

\subsection*{You'll want to know: \sk{राजते}}
\textbf{\sk{राजते}} — \emph{rājate} — \textquotedblleft{}gleams; shines brightly\textquotedblright{}.

\begin{quote}
\sk{अमावास्यायाः} \sk{अनन्तरं} \sk{चन्द्रः} \sk{पुनः} \sk{राजते।}
\end{quote}

\begin{quote}
\emph{amāvāsyāyāḥ anantaraṃ candraḥ punaḥ rājate.}
\end{quote}

\begin{quote}
After the new moon the moon gleams again.
\end{quote}

Past (it …): \sk{चन्द्रः} \sk{पुनः} \sk{राजितवान्}, \emph{candraḥ punaḥ rājitavān}.

\begin{cousinweb}
\textbf{\sk{राजते}} (\emph{rājate}) comes from the root √rāj, \textquotedblleft{}to shine, to be splendid\textquotedblright{}, which MW also gives as to reign, and thinks was perhaps two roots at first; MW gives the participle \emph{rājita}, \textquotedblleft{}resplendent\textquotedblright{} (MW).
\end{cousinweb}

\begin{grammarlens}[title={What you've built}]
One more word for getting around.
\end{grammarlens}

\subsection*{Your turn}
\begin{itemize}
\raggedright
  \item \emph{Say it:} \emph{rājate}
  \item \emph{Say it:} \emph{rājate}, once more
  \item \emph{Say it:} \emph{amāvāsyāyāḥ anantaraṃ candraḥ punaḥ rājate.} — After the new moon the moon gleams again.
  \item \emph{Recall:} say the Sanskrit for the bright fortnight, then the Sanskrit for the dark fortnight, then say \emph{rājate} again
\end{itemize}

\begin{tcolorbox}[breakable,colback=teal!4,colframe=teal!35!black,title={Before you move on}]
What does \sk{राजते} mean? (\textquotedblleft{}Gleams; shines brightly\textquotedblright{}.) Say it once more.
\end{tcolorbox}
